% ===========================================================================
%  comune/preambolo.tex — impostazioni comuni ai due documenti
% ===========================================================================
%
%  Lo includono sia guida-sviluppatore.tex sia manuale-utente.tex. Ciò che
%  vale per uno solo dei due NON va qui: va nel suo file principale, altrimenti
%  fra sei mesi nessuno sa più perché una cosa è dove è.
%
%  Motore: pdfLaTeX. Non lualatex, non xelatex, e la scelta è deliberata —
%  questo preambolo usa solo pacchetti di una TeX Live di base, così chi riceve
%  la cartella compila senza installare nulla. Nessun font di sistema, nessun
%  --shell-escape, nessuna dipendenza da Python.
% ===========================================================================

% --- lingua e codifica ------------------------------------------------------
\usepackage[T1]{fontenc}
\usepackage[utf8]{inputenc}

%  L'italiano in due modi, perché non tutte le installazioni hanno il pacchetto
%  babel-italian (le immagini Docker di TeX Live «base» tipicamente no, e una
%  documentazione che non compila lì è inutile metà delle volte che serve).
%    - c'è italian.ldf  -> si usa quello, è la via normale;
%    - non c'è          -> babel >= 3.75 costruisce la lingua dal proprio file
%                          locale/it/babel-it.ini, che è sempre presente.
%  In entrambi i casi si ottengono «Capitolo», «Indice», «Figura», «Tabella».
%  La sillabazione italiana è un'altra cosa: sta nel FORMATO, non nel pacchetto,
%  e se manca si vede subito (parole spezzate all'inglese). Il controllo è in
%  fondo a questo file.
\IfFileExists{italian.ldf}
  {\usepackage[italian]{babel}}
  {\usepackage[italian, provide = *]{babel}}

\usepackage{csquotes}

% --- caratteri --------------------------------------------------------------
%  Palatino per il testo e per la matematica (mathpazo): la documentazione ha
%  formule — vol targeting, restringimento della covarianza — e un testo con la
%  matematica di un'altra famiglia si legge come due documenti incollati.
%  Helvetica ridotta per i titoli, monospazio di Latin Modern per il codice.
\usepackage{mathpazo}
\usepackage[scaled=0.90]{helvet}
\usepackage{lmodern}
\renewcommand{\ttdefault}{lmtt}

\usepackage{microtype}

% --- composizione del paragrafo ---------------------------------------------
%  Un documento pieno di nomi di file e di simboli in monospazio ha molte
%  parole che non si possono spezzare: `diagnose_regime_signals.py` e' lungo
%  quattro centimetri e non ha punti di sillabazione. Con i valori standard
%  TeX preferisce far sbordare la riga nel margine piuttosto che allargare
%  gli spazi, e il risultato si vede.
%
%  Questi tre valori invertono la preferenza: meglio una riga un po' larga di
%  spazi che una riga che esce dalla giustezza. `hbadness` alza la soglia di
%  segnalazione di conseguenza, cosi' il log continua a riportare i problemi
%  veri e non le righe lente che abbiamo appena autorizzato.
\tolerance = 1500
\emergencystretch = 3em
\hbadness = 2000

% --- impaginazione ----------------------------------------------------------
\usepackage[
  a4paper,
  inner = 28mm, outer = 24mm, top = 26mm, bottom = 28mm,
  headsep = 7mm, footskip = 12mm,
]{geometry}

\usepackage{emptypage}      % le pagine lasciate bianche restano davvero bianche

% Niente pagine bianche. I due documenti si leggono a schermo: la pagina
% vuota che il libro inserisce per far cominciare capitoli, parti e
% frontespizio sulla pagina di destra lì è solo un buco. Con l'opzione
% openany della classe i capitoli e le parti cominciano sulla pagina
% seguente; qui lo stesso vale per \frontmatter, \mainmatter e i
% \cleardoublepage scritti a mano (frontespizi), che la classe book
% eseguirebbe comunque. Il layout resta a due facciate (margini interno ed
% esterno, intestazioni alternate).
\let\cleardoublepage\clearpage
\usepackage{setspace}
\setstretch{1.06}

% --- colori, codice, figure -------------------------------------------------
\usepackage{etoolbox}
% ===========================================================================
%  comune/colori.tex — la tavolozza
% ===========================================================================
%
%  Sono i colori del PRODOTTO, non colori scelti per la documentazione: gli
%  stessi valori che stanno in dashboard/static/css/ui-v2.css e nella landing.
%  Documentare un prodotto con i colori di qualcun altro costringe chi legge a
%  tradurre due volte — una figura verde acqua qui e un pulsante verde acqua
%  nell'app sono la stessa cosa, e si deve vedere.
%
%  Chi cambia la tavolozza del prodotto cambia QUESTO file e nient'altro:
%  nessun altro file della documentazione nomina un colore per esteso.
% ===========================================================================
\RequirePackage{xcolor}

% --- i colori di base della landing -----------------------------------------
\definecolor{lyraInchiostro}{HTML}{111318}  % testo, titoli
\definecolor{lyraPannello}{HTML}{F5F5F2}    % fondo dei riquadri
\definecolor{lyraBordo}{HTML}{ECECEA}       % filetti e bordi
\definecolor{lyraAccento}{HTML}{2BB59A}     % verde acqua: l'unico accento
\definecolor{lyraPositivo}{HTML}{0F7A62}    % numeri in utile
\definecolor{lyraNegativo}{HTML}{9A3412}    % numeri in perdita

% --- i pastelli, usati per distinguere gli strati nelle figure --------------
\definecolor{lyraCielo}{HTML}{D9F1F7}
\definecolor{lyraMenta}{HTML}{DDF3E4}
\definecolor{lyraPesca}{HTML}{FFE4D4}
\definecolor{lyraLavanda}{HTML}{E6DEFF}
\definecolor{lyraCrema}{HTML}{FFF1C9}

% --- derivati, per bordi e testi sopra i pastelli ---------------------------
%  Si ricavano dai precedenti invece di essere scritti a mano: cambiando un
%  pastello, il suo bordo lo segue.
\colorlet{lyraCieloBordo}{lyraCielo!55!lyraInchiostro}
\colorlet{lyraMentaBordo}{lyraMenta!55!lyraInchiostro}
\colorlet{lyraPescaBordo}{lyraPesca!55!lyraInchiostro}
\colorlet{lyraLavandaBordo}{lyraLavanda!55!lyraInchiostro}
\colorlet{lyraCremaBordo}{lyraCrema!55!lyraInchiostro}

\colorlet{lyraGrigio}{lyraInchiostro!55}
\colorlet{lyraGrigioChiaro}{lyraInchiostro!28}
\colorlet{lyraAccentoScuro}{lyraAccento!70!lyraInchiostro}
\colorlet{lyraAccentoTenue}{lyraAccento!14!white}

% --- i colori del codice ----------------------------------------------------
\colorlet{codFondo}{lyraPannello}
\colorlet{codBordo}{lyraBordo}
\colorlet{codCommento}{lyraInchiostro!45}
\colorlet{codParola}{lyraAccentoScuro}
\colorlet{codStringa}{lyraNegativo!85}
\colorlet{codNumero}{lyraInchiostro!70}
\colorlet{codRiga}{lyraInchiostro!35}

% ===========================================================================
%  comune/stile-tikz.tex — il vocabolario grafico delle figure
% ===========================================================================
%
%  Ogni figura del progetto disegna con QUESTI stili e non con parametri
%  propri. La ragione e' leggibilita': se in una figura il riquadro di un
%  modulo e' azzurro e in quella dopo e' dello stesso azzurro un archivio su
%  disco, chi legge impara due volte e si fida zero. Qui un colore significa
%  una cosa sola in tutto il documento:
%
%     cielo     il MOTORE di calcolo (Python, core/)
%     menta     l'INTERFACCIA (Flask, dashboard/)
%     pesca     il SITO (Next.js) e tutto cio' che sta in rete
%     lavanda   i dati PERSISTENTI (disco, Postgres)
%     crema     cio' che viene da FUORI (broker, Stripe, provider dati)
%     bianco    strati astratti (porte, contratti, protocolli)
%
%  Questo file e' incluso sia dai documenti sia dalle figure compilate da sole
%  (figure/tikz/*.tex con la classe standalone), percio' non deve dipendere da
%  niente che esista solo dentro un documento.
% ===========================================================================
\RequirePackage{tikz}
\usetikzlibrary{
  arrows.meta, positioning, fit, backgrounds, calc, shapes.geometric,
  shapes.multipart, shapes.symbols, decorations.pathreplacing,
  decorations.pathmorphing, patterns, shadows.blur, matrix, chains, scopes
}

\tikzset{
  % --- impostazioni generali ----------------------------------------------
  every picture/.style = {
    line width = 0.7pt,
    font = \sffamily\small,
  },
  %
  % --- RIQUADRI: un modulo, un componente, una responsabilita' -------------
  riquadro/.style = {
    rectangle, rounded corners = 3pt, draw,
    minimum height = 9.5mm, minimum width = 30mm,
    inner xsep = 6pt, inner ysep = 5pt,
    align = center, font = \sffamily\small,
    text = lyraInchiostro,
  },
  motore/.style    = {riquadro, fill = lyraCielo,   draw = lyraCieloBordo},
  interfaccia/.style = {riquadro, fill = lyraMenta,  draw = lyraMentaBordo},
  sito/.style      = {riquadro, fill = lyraPesca,   draw = lyraPescaBordo},
  persistenza/.style = {riquadro, fill = lyraLavanda, draw = lyraLavandaBordo},
  esterno/.style   = {riquadro, fill = lyraCrema,   draw = lyraCremaBordo},
  astratto/.style  = {riquadro, fill = white, draw = lyraGrigio, dashed},
  %
  % Il riquadro piccolo: per gli elenchi fitti dove il nome basta.
  minuto/.style = {riquadro, minimum height = 7mm, minimum width = 22mm,
                   font = \sffamily\scriptsize, inner xsep = 4pt},
  %
  % --- FORME PARTICOLARI ---------------------------------------------------
  archivio/.style = {
    cylinder, shape border rotate = 90, aspect = 0.22, draw = lyraLavandaBordo,
    fill = lyraLavanda, minimum height = 13mm, minimum width = 22mm,
    inner xsep = 0pt, inner ysep = 4pt, align = center,
    font = \sffamily\scriptsize, text = lyraInchiostro,
  },
  decisione/.style = {
    diamond, draw = lyraInchiostro!60, fill = white, aspect = 1.9,
    inner xsep = 2pt, inner ysep = 0pt, align = center,
    font = \sffamily\scriptsize, text = lyraInchiostro,
  },
  persona/.style = {
    rectangle, rounded corners = 10pt, draw = lyraInchiostro!70,
    fill = white, minimum height = 9mm, minimum width = 24mm,
    align = center, font = \sffamily\small\itshape,
  },
  %
  % --- STRATI: il contenitore che raggruppa -------------------------------
  strato/.style = {
    rectangle, rounded corners = 5pt, draw = lyraGrigioChiaro,
    fill = lyraPannello!55!white, inner sep = 9pt,
  },
  stratoTitolo/.style = {
    font = \sffamily\scriptsize\bfseries, text = lyraInchiostro!65,
    anchor = north west,
  },
  %
  % --- FRECCE --------------------------------------------------------------
  %  Tre tipi, e la differenza va vista senza leggere la legenda:
  %    flusso      pieno    -> i dati passano di qui
  %    dipendenza  tratteg. -> questo conosce quello (ma nulla scorre)
  %    ritorno     grigio   -> la risposta che torna indietro
  flusso/.style = {
    -{Stealth[length = 2.4mm, width = 1.8mm]},
    draw = lyraInchiostro!75, line width = 0.8pt,
  },
  dipendenza/.style = {
    -{Stealth[length = 2.2mm, width = 1.6mm]}, dashed, dash pattern = on 2.4pt off 2pt,
    draw = lyraInchiostro!55, line width = 0.6pt,
  },
  ritorno/.style = {
    -{Stealth[length = 2.2mm, width = 1.6mm]},
    draw = lyraGrigio, line width = 0.6pt,
  },
  accento/.style = {
    -{Stealth[length = 2.6mm, width = 2mm]},
    draw = lyraAccento, line width = 1pt,
  },
  bloccato/.style = {
    -{Bar[width = 2.4mm]}, draw = lyraNegativo, line width = 0.9pt,
  },
  %
  % --- ETICHETTE SULLE FRECCE ---------------------------------------------
  etichetta/.style = {
    font = \sffamily\scriptsize, text = lyraInchiostro!70,
    fill = white, inner xsep = 2.5pt, inner ysep = 1pt,
  },
  etichettaAccento/.style = {etichetta, text = lyraAccentoScuro},
  %
  % --- NOTE A MARGINE ------------------------------------------------------
  nota/.style = {
    font = \sffamily\scriptsize\itshape, text = lyraInchiostro!60,
    align = left, inner sep = 2pt,
  },
  notaAccento/.style = {nota, text = lyraAccentoScuro, font = \sffamily\scriptsize},
  %
  % --- SEGNAPOSTO NUMERICI (1) (2) (3) sulle figure a passi ---------------
  passo/.style = {
    circle, draw = lyraAccento, fill = white, text = lyraAccentoScuro,
    font = \sffamily\scriptsize\bfseries, inner sep = 1.2pt,
    minimum size = 4.6mm,
  },
}

% --- la cornice con il titolo, usata dalle figure a strati -------------------
%  #1 opzioni  #2 nodi da contenere  #3 titolo
%  Il titolo sta SOPRA la cornice, non dentro: dentro finisce addosso al primo
%  nodo ogni volta che il nodo è vicino al bordo, e lo si scopre solo guardando
%  il PDF. Il nome del nodo è contato, così due strati nella stessa figura non
%  si sovrascrivono a vicenda.
\newcounter{lyrastratoN}
\newcommand{\lyrastrato}[3][]{%
  \stepcounter{lyrastratoN}%
  \begin{scope}[on background layer]
    \node[strato, fit = #2, #1] (cornice\thelyrastratoN) {};
    \node[stratoTitolo, anchor = south west]
      at ([shift={(1pt,1.5pt)}]cornice\thelyrastratoN.north west) {#3};
  \end{scope}}

% --- legenda compatta, da mettere sotto una figura ---------------------------
%  Si usa solo quando la figura introduce un colore nuovo: ripeterla sotto ogni
%  figura la fa diventare rumore che nessuno legge piu'.
\newcommand{\lyralegenda}[1]{%
  \begin{tikzpicture}[baseline]
    \matrix[matrix of nodes, nodes = {anchor = west, font = \sffamily\scriptsize},
            column sep = 3mm, row sep = 0.6mm] {#1\\};
  \end{tikzpicture}}

% ---------------------------------------------------------------------------
%  DUE TRAPPOLE, scritte qui perché costano mezz'ora ciascuna e l'errore che
%  producono indica sempre il punto sbagliato.
%
%  1. `\\` dentro un gruppo del contenuto di un nodo.
%        {testo\\{\scriptsize una riga\\e un'altra}}   <-- SI ROMPE
%        {testo\\{\scriptsize una riga}\\{\scriptsize e un'altra}}  <-- giusto
%     TikZ raccoglie le righe guardando i `\\` al primo livello: quello dentro
%     le graffe è il `\\` di LaTeX, che in un box orizzontale non ha senso.
%     L'errore che ne esce è «\tikzscope@linewidth non definito», segnalato sul
%     nodo DOPO quello colpevole.
%
%  2. Una tcolorbox dentro un nodo. Vale per \cod, \pulsante, \campo: sono
%     tcbox, cioè figure TikZ annidate, e lasciano la figura esterna in uno
%     stato incoerente. Stesso errore, stessa confusione. Nelle figure si usa
%     \codT (qui sotto).
% ---------------------------------------------------------------------------

% --- il codice DENTRO una figura --------------------------------------------
%  Nel testo un simbolo si scrive \cod{...}, che è un \tcbox: una tcolorbox,
%  cioè una tikzpicture. Infilarla dentro il contenuto di un nodo TikZ rompe lo
%  stato della figura che la contiene, e l'errore salta fuori sul nodo
%  SUCCESSIVO (`\tikzscope@linewidth` non definito) — mezz'ora buttata a
%  cercarlo nel posto sbagliato. Dentro le figure si usa questo, che è solo un
%  cambio di carattere e non apre niente.
%  \upshape serve davvero: dentro un nodo `nota`, che è in corsivo, un
%  monospazio corsivo si legge male e sembra un'altra cosa.
\newcommand{\codT}[1]{{\ttfamily\upshape\small\color{lyraInchiostro!85}#1}}

\newcommand{\lyravoce}[2]{%
  \tikz[baseline=-0.55ex]\node[#1, minimum width = 4.2mm, minimum height = 2.6mm,
        inner sep = 0pt, rounded corners = 1.4pt]{};~#2}

% ===========================================================================
%  comune/stile-codice.tex — come si mostra il codice
% ===========================================================================
%
%  PERCHE' listings E NON minted. minted colora meglio, ma per farlo chiama
%  Pygments durante la compilazione: serve Python installato e il permesso
%  --shell-escape, che molti editor LaTeX negano per difetto e che le
%  compilazioni automatiche rifiutano. Una documentazione che non compila sulla
%  macchina di chi la legge non e' documentazione. listings sta dentro LaTeX e
%  compila ovunque, e la differenza di resa non vale quel prezzo.
%
%  Tre modi di mostrare codice, e vanno usati per cose diverse:
%
%    \begin{codice}[python]{didascalia}{percorso/file.py}  blocco con cornice
%    \inputcodice{python}{didascalia}{file.tex}{percorso}  da file esterno
%    \cod{nome_simbolo}                                     dentro una frase
%
%  Il PERCORSO non e' decorazione: chi legge la documentazione ha il sorgente
%  aperto accanto, e un frammento senza il suo file costringe a cercarlo.
% ===========================================================================
\RequirePackage{listings}
\RequirePackage{tcolorbox}
\tcbuselibrary{listings, skins, breakable}

% --- stile comune ----------------------------------------------------------
\lstdefinestyle{lyraBase}{
  basicstyle      = \ttfamily\footnotesize\color{lyraInchiostro},
  keywordstyle    = \color{codParola}\bfseries,
  commentstyle    = \color{codCommento}\itshape,
  stringstyle     = \color{codStringa},
  numberstyle     = \ttfamily\scriptsize\color{codRiga},
  numbers         = left,
  numbersep       = 9pt,
  stepnumber      = 1,
  showstringspaces= false,
  breaklines      = true,
  breakatwhitespace = false,
  % La freccia dice che la riga continua: senza, una riga spezzata si legge
  % come due istruzioni, ed e' il tipo di errore che si propaga nel copia-incolla.
  prebreak        = \raisebox{0ex}[0ex][0ex]{\color{lyraAccento}\ensuremath{\hookleftarrow}},
  tabsize         = 4,
  columns         = fullflexible,
  keepspaces      = true,
  upquote         = true,
  extendedchars   = true,
  literate        =
    {à}{{\`a}}1 {è}{{\`e}}1 {é}{{\'e}}1 {ì}{{\`\i}}1 {ò}{{\`o}}1 {ù}{{\`u}}1
    {À}{{\`A}}1 {È}{{\`E}}1 {Ì}{{\`I}}1 {Ò}{{\`O}}1 {Ù}{{\`U}}1
    {€}{{\texteuro}}1 {°}{{\textdegree}}1 {·}{{\textperiodcentered}}1
    {—}{{---}}1 {–}{{--}}1 {«}{{\guillemotleft}}1 {»}{{\guillemotright}}1
    {“}{{\textquotedblleft}}1 {”}{{\textquotedblright}}1
    {’}{{\textquoteright}}1 {‘}{{\textquoteleft}}1
    {→}{{$\rightarrow$}}1 {←}{{$\leftarrow$}}1 {✓}{{$\checkmark$}}1,
}

% ---------------------------------------------------------------------------
%  REGOLA: dentro un blocco di codice si scrive SOLO ASCII.
%
%  Niente lettere accentate, niente caratteri di disegno, niente frecce.
%  `literate` non basta, e nemmeno il pacchetto listingsutf8: tcblisting
%  scrive il contenuto in un file intermedio e lo rilegge, quindi i byte
%  arrivano al kernel prima che listings possa intercettarli. L'errore e'
%  «Invalid UTF-8 byte sequence», segnalato nel file .listing e senza dire
%  quale blocco l'ha causato.
%
%  Quando un valore del sorgente contiene un accento, si translittera e si
%  mette un commento accanto. Gli alberi di cartelle si disegnano in ASCII
%  (|-- e `--), oppure in TikZ, dove non passano da listings.
%  `make controlla` segnala le violazioni prima che fermino la compilazione.
% ---------------------------------------------------------------------------

% --- i linguaggi che compaiono nel progetto ---------------------------------
\lstdefinelanguage{JavaScriptLyra}{
  keywords = {break, case, catch, class, const, continue, debugger, default,
    delete, do, else, export, extends, finally, for, function, if, import, in,
    instanceof, let, new, of, return, super, switch, this, throw, try, typeof,
    var, void, while, with, yield, async, await, static, get, set},
  keywordstyle = \color{codParola}\bfseries,
  ndkeywords = {true, false, null, undefined, document, window, console,
    fetch, JSON, Promise, Array, Object, Math},
  ndkeywordstyle = \color{lyraInchiostro!75}\bfseries,
  sensitive = true,
  comment = [l]{//},
  morecomment = [s]{/*}{*/},
  morestring = [b]',
  morestring = [b]",
  morestring = [b]`,
}

\lstdefinelanguage{TypeScriptLyra}{
  morekeywords = [1]{break, case, catch, class, const, continue, default,
    delete, do, else, enum, export, extends, finally, for, function, if,
    implements, import, in, instanceof, interface, let, new, of, private,
    protected, public, readonly, return, super, switch, this, throw, try,
    type, typeof, var, void, while, yield, async, await, as, satisfies},
  morekeywords = [2]{string, number, boolean, any, unknown, never, Promise,
    Record, Array, Partial, Required},
  keywordstyle = [1]\color{codParola}\bfseries,
  keywordstyle = [2]\color{lyraInchiostro!70},
  sensitive = true,
  comment = [l]{//},
  morecomment = [s]{/*}{*/},
  morestring = [b]',
  morestring = [b]",
  morestring = [b]`,
}

\lstdefinelanguage{Console}{
  basicstyle = \ttfamily\footnotesize,
  morecomment = [l]{\#},
  sensitive = false,
}

% --- l'ambiente principale --------------------------------------------------
%  #1 linguaggio, con il nome CORTO (python, js, ts, sql, json, yaml, html,
%     css, console): \lstlingua lo traduce nel nome che listings conosce, così
%     chi scrive non deve ricordare che JSON non esiste e che gli somiglia Java
%  #2 didascalia
%  #3 percorso del file nel progetto (vuoto = nessuno)
\newtcblisting{codiceInterno}[3]{
  listing only,
  % Sia lo stile sia la lingua devono stare QUI dentro: `listing options`
  % sostituisce `listing style` invece di aggiungersi, e il sintomo e' che
  % meta' delle impostazioni (breaklines, numeri di riga) sparisce senza
  % dirlo. Le righe lunghe escono dal riquadro, ed e' l'unico segno.
  listing options = {style = lyraBase, language = #1},
  breakable,
  enhanced,
  colback   = codFondo,
  colframe  = codBordo,
  boxrule   = 0.5pt,
  arc       = 3pt,
  left      = 3pt, right = 3pt, top = 4pt, bottom = 4pt,
  toptitle  = 2pt, bottomtitle = 2pt,
  title     = {\footnotesize\sffamily\bfseries\color{lyraInchiostro}#2%
               \ifstrempty{#3}{}{\hfill\normalfont\ttfamily\scriptsize
                 \color{lyraInchiostro!60}#3}},
  colbacktitle = lyraPannello!60!white,
  coltitle  = lyraInchiostro,
  fonttitle = \footnotesize,
  titlerule = 0pt,
  before skip = 10pt, after skip = 10pt,
}

% Nomi corti per i linguaggi, cosi' chi scrive non deve ricordare i nomi
% interni di listings ne' che JSON e YAML non esistono.
%
% Si usa \lstalias e non una macro che traduce: l'opzione `language` di
% listings viene letta senza espandere, quindi una macro li' dentro produce
% «Missing \endcsname» alla fine del blocco di codice — errore che non dice
% niente di cio' che e' successo.
\lstalias{console}{Console}
\lstalias{python}{Python}
\lstalias{js}{JavaScriptLyra}
\lstalias{ts}{TypeScriptLyra}
\lstalias{sql}{SQL}
\lstalias{json}{Java}      % JSON non esiste in listings: Java gli somiglia
\lstalias{yaml}{Python}    % idem per YAML: commenti con #, due punti, indentazione
\lstalias{html}{HTML}
\lstalias{css}{HTML}
\lstalias{testo}{}

% --- codice da file esterno -------------------------------------------------
%  #1 linguaggio  #2 didascalia  #3 percorso del file su disco  #4 percorso da
%  mostrare. Si usa per gli estratti lunghi tenuti fuori dal testo.
\newcommand{\inputcodice}[4]{%
  \begin{tcolorbox}[
    breakable, enhanced,
    colback = codFondo, colframe = codBordo, boxrule = 0.5pt, arc = 3pt,
    left = 3pt, right = 3pt, top = 4pt, bottom = 4pt,
    title = {\footnotesize\sffamily\bfseries\color{lyraInchiostro}#2%
             \hfill\normalfont\ttfamily\scriptsize\color{lyraInchiostro!60}#4},
    colbacktitle = lyraPannello!60!white, coltitle = lyraInchiostro,
    fonttitle = \footnotesize, titlerule = 0pt]
    \lstinputlisting[style=lyraBase, language=#1]{#3}
  \end{tcolorbox}}

% --- codice dentro una frase ------------------------------------------------
%  Fondo tenue invece del grassetto: dentro un paragrafo il grassetto chiama
%  l'occhio e spezza la lettura, mentre un nome di funzione va riconosciuto
%  senza essere gridato.
%  Le due penalita' attorno alla scatola non sono decorazione. Un \tcbox e'
%  una scatola indivisibile: se capita dopo una parentesi o un trattino --
%  cioe' senza spazio davanti -- TeX non ha dove spezzare la riga e la lascia
%  sbordare nel margine. Sono state la causa di quasi tutte le righe troppo
%  lunghe della prima stesura. Con queste, la scatola puo' sempre andare a
%  capo per conto suo.
%
%  Resta un limite: una scatola piu' larga della giustezza sborda comunque.
%  Per un percorso lungo si usa \file{} (testo normale, si spezza) oppure lo
%  si mette in un blocco di codice.
\DeclareRobustCommand{\cod}[1]{%
  \penalty50
  \tcbox[on line, boxsep = 0.6pt, left = 1.4pt, right = 1.4pt, top = 0pt,
         bottom = 0pt, arc = 2pt, boxrule = 0pt,
         colback = lyraPannello, colframe = lyraPannello]
        {\ttfamily\small #1}%
  \penalty0 }

% Percorso di file nel testo.
%  Usa lo stesso ciclo di \percorsoHTTP (definito in comune/macro.tex): un nome
%  come `diagnose_regime_signals.py' e' lungo quattro centimetri e senza punti
%  di sillabazione, quindi in una colonna stretta sborda. Con i punti di
%  interruzione dopo / _ . e - va a capo dove si legge comunque.
%  Il ciclo richiede ASCII e niente spazi: per il resto c'e' \cod{}.
\DeclareRobustCommand{\file}[1]{{\color{lyraInchiostro!78}\percorsoHTTP{#1}}}

% ===========================================================================
%  Un COMMENTO del sorgente citato per esteso
% ===========================================================================
%  Diversa da `codiceInterno`, e la differenza non e' estetica. Molti commenti
%  di questo progetto sono prose italiane di dieci righe che spiegano una
%  decisione: sono la parte piu' utile da citare, e dentro un blocco di
%  listings non ci stanno — listings vuole ASCII puro (vedi sopra), e
%  translitterare «perche'» e le lettere accentate significa citare qualcosa
%  che nel sorgente non c'e'.
%
%  Qui il testo e' testo LaTeX normale: accenti, virgolette e trattini lunghi
%  si compongono come nel resto del documento, e la citazione resta fedele.
%
%    \begin{commento}{percorso/del/file.py:120}
%    Il testo del commento, a capo come nel sorgente.
%    \end{commento}
%
%  Il corsivo e il filetto a sinistra dicono «questo l'ha scritto chi ha
%  scritto il codice, non chi ha scritto la guida».
\newtcolorbox{commento}[1]{
  enhanced, breakable,
  colback   = codFondo, colframe = lyraInchiostro!30,
  boxrule   = 0pt, leftrule = 2.2pt, arc = 2pt,
  left = 9pt, right = 9pt, top = 6pt, bottom = 6pt,
  fontupper = \small\itshape,
  title     = {\normalfont\ttfamily\scriptsize\color{lyraInchiostro!60}#1},
  colbacktitle = codFondo, coltitle = lyraInchiostro,
  titlerule = 0pt, toptitle = 2pt, bottomtitle = 0pt,
  before skip = 10pt, after skip = 10pt,
}


% --- numeri -----------------------------------------------------------------
%  In italiano il separatore delle migliaia sarebbe il punto, ma in un
%  documento pieno di valori numerici (soglie, quantili, prezzi) il punto
%  decimale inglese del codice e il punto delle migliaia del testo si
%  confondono. Spazio fine: non si confonde con niente.
\usepackage{siunitx}
\sisetup{group-separator = {\,}, group-minimum-digits = 4, detect-all}

% --- tabelle ----------------------------------------------------------------
\usepackage{booktabs}
\usepackage{longtable}
\usepackage{multirow}
\usepackage{array}
\usepackage{tabularx}
\usepackage{colortbl}

\newcolumntype{L}[1]{>{\raggedright\arraybackslash}p{#1}}
\newcolumntype{C}[1]{>{\centering\arraybackslash}p{#1}}
\newcolumntype{R}[1]{>{\raggedleft\arraybackslash}p{#1}}
\newcolumntype{M}{>{\ttfamily\footnotesize}l}

% Righe alternate tenui: su una tabella di venti righe l'occhio perde la riga,
% e un colore appena percepibile costa meno di una griglia piena.
\newcommand{\rigaTenue}{\rowcolor{lyraPannello!55!white}}

% --- elenchi ----------------------------------------------------------------
\usepackage{enumitem}
\setlist{itemsep = 1pt, parsep = 2pt, topsep = 4pt}
\setlist[itemize]{label = \textcolor{lyraAccento}{\textbullet}}
\setlist[enumerate]{label = \textcolor{lyraAccentoScuro}{\arabic*.}}

% --- figure -----------------------------------------------------------------
\usepackage{graphicx}
\usepackage{float}
\usepackage[font = small, labelfont = {bf, sf}, textfont = sf,
            labelsep = period, skip = 7pt]{caption}
\usepackage{subcaption}
\usepackage{adjustbox}
\graphicspath{{figure/}{figure/tikz/}{figure/schermate/}}

% --- intestazioni e piè di pagina -------------------------------------------
\usepackage{fancyhdr}
\usepackage{titlesec}
\usepackage{titletoc}

% --- riferimenti ------------------------------------------------------------
\usepackage[
  colorlinks = true,
  linkcolor  = lyraAccentoScuro,
  citecolor  = lyraAccentoScuro,
  urlcolor   = lyraAccentoScuro,
  filecolor  = lyraAccentoScuro,
  breaklinks = true,
  bookmarksnumbered = true,
  bookmarksopen = true,
  bookmarksopenlevel = 1,
]{hyperref}
\usepackage{url}
\urlstyle{sf}

% --- indice analitico -------------------------------------------------------
\usepackage{imakeidx}

% ===========================================================================
%  Stile dei titoli
% ===========================================================================
%  I capitoli si aprono con il numero grande a sinistra e un filetto verde
%  acqua: è il segno che si gira pagina in una sezione nuova, e lo si vede
%  sfogliando senza leggere.
\titleformat{\chapter}[display]
  {\sffamily\LARGE\bfseries\color{lyraInchiostro}}
  {\raggedright\normalfont\sffamily\fontsize{52}{52}\selectfont
   \color{lyraAccento!35!lyraPannello}\thechapter}
  {-6pt}
  {\raggedright}
  [\vspace{4pt}{\color{lyraAccento}\rule{\textwidth}{1.2pt}}]
\titlespacing*{\chapter}{0pt}{-24pt}{26pt}

\titleformat{\section}
  {\sffamily\large\bfseries\color{lyraInchiostro}}
  {\color{lyraAccentoScuro}\thesection}{0.7em}{}
\titlespacing*{\section}{0pt}{18pt}{7pt}

\titleformat{\subsection}
  {\sffamily\normalsize\bfseries\color{lyraInchiostro!88}}
  {\color{lyraAccentoScuro}\thesubsection}{0.6em}{}
\titlespacing*{\subsection}{0pt}{13pt}{4pt}

\titleformat{\subsubsection}[runin]
  {\sffamily\small\bfseries\color{lyraInchiostro!80}}
  {}{0pt}{}[\;---\;]
\titlespacing*{\subsubsection}{0pt}{9pt}{0.5em}

\setcounter{secnumdepth}{3}
\setcounter{tocdepth}{2}

% ===========================================================================
%  Indice
% ===========================================================================
%  L'indice è composto con titletoc, non con tocloft: i due pacchetti si
%  contendono le stesse macro e caricarli insieme rompe l'indice in modi che
%  si vedono solo a documento finito.
\titlecontents{chapter}[0pt]
  {\addvspace{9pt}\sffamily\bfseries}
  {\color{lyraAccentoScuro}\contentslabel{1.6em}}
  {}
  {\color{lyraInchiostro!55}\titlerule*[0.6pc]{.}\contentspage}
\titlecontents{section}[1.9em]
  {\sffamily\small}
  {\color{lyraInchiostro!70}\contentslabel{2.3em}}
  {}
  {\color{lyraInchiostro!45}\titlerule*[0.6pc]{.}\contentspage}

% ===========================================================================
%  I riquadri del testo
% ===========================================================================
%  Quattro soli, e ciascuno ha un lavoro preciso. Un documento con nove tipi di
%  riquadro è un documento in cui nessun riquadro significa qualcosa.
%
%    \begin{nota}{titolo}        un chiarimento utile, saltabile
%    \begin{attenzione}{titolo}  qui si rompe qualcosa
%    \begin{perche}{titolo}      perché è così e non altrimenti
%    \begin{pratica}{titolo}     cosa fare, passo per passo
% ===========================================================================
\newtcolorbox{nota}[1]{
  enhanced, breakable,
  colback = lyraCielo!35!white, colframe = lyraCieloBordo,
  boxrule = 0pt, leftrule = 2.2pt, arc = 2pt,
  left = 8pt, right = 8pt, top = 6pt, bottom = 6pt,
  fonttitle = \sffamily\bfseries\small,
  coltitle = lyraInchiostro,
  title = {#1}, titlerule = 0pt, colbacktitle = lyraCielo!35!white,
  toptitle = 1pt, bottomtitle = 1pt,
  before skip = 10pt, after skip = 10pt,
}

\newtcolorbox{attenzione}[1]{
  enhanced, breakable,
  colback = lyraPesca!40!white, colframe = lyraNegativo!70,
  boxrule = 0pt, leftrule = 2.2pt, arc = 2pt,
  left = 8pt, right = 8pt, top = 6pt, bottom = 6pt,
  fonttitle = \sffamily\bfseries\small,
  coltitle = lyraNegativo,
  title = {#1}, titlerule = 0pt, colbacktitle = lyraPesca!40!white,
  toptitle = 1pt, bottomtitle = 1pt,
  before skip = 10pt, after skip = 10pt,
}

\newtcolorbox{perche}[1]{
  enhanced, breakable,
  colback = lyraPannello, colframe = lyraGrigioChiaro,
  boxrule = 0pt, leftrule = 2.2pt, arc = 2pt,
  left = 8pt, right = 8pt, top = 6pt, bottom = 6pt,
  fonttitle = \sffamily\bfseries\small,
  coltitle = lyraInchiostro!75,
  title = {#1}, titlerule = 0pt, colbacktitle = lyraPannello,
  toptitle = 1pt, bottomtitle = 1pt,
  before skip = 10pt, after skip = 10pt,
}

\newtcolorbox{pratica}[1]{
  enhanced, breakable,
  colback = lyraMenta!45!white, colframe = lyraAccento,
  boxrule = 0pt, leftrule = 2.2pt, arc = 2pt,
  left = 8pt, right = 8pt, top = 6pt, bottom = 6pt,
  fonttitle = \sffamily\bfseries\small,
  coltitle = lyraAccentoScuro,
  title = {#1}, titlerule = 0pt, colbacktitle = lyraMenta!45!white,
  toptitle = 1pt, bottomtitle = 1pt,
  before skip = 10pt, after skip = 10pt,
}

% ===========================================================================
%  Macro del progetto
% ===========================================================================
% ===========================================================================
%  comune/macro.tex — le abbreviazioni del progetto
% ===========================================================================
%
%  Servono a scrivere una cosa in un posto solo. Il nome di un pacchetto
%  Python compare duecento volte nella guida: se un giorno `approaches/`
%  diventa `overlays/`, si cambia qui e non duecento volte a mano — e
%  soprattutto non si finisce con metà documento che dice una cosa e metà
%  un'altra, che è il modo in cui la documentazione smette di essere creduta.
% ===========================================================================

% --- la lingua del documento -----------------------------------------------
%  \iflyraInglese la definisce comune/preambolo.tex. Qui la si crea se manca,
%  perché `make figure' include questo file senza il preambolo. Il nome del
%  condizionale si scrive con \csname: scritto per esteso dentro un ramo
%  saltato, TeX lo conterebbe come un \if aperto e sbaglierebbe i \fi.
\ifcsname iflyraInglese\endcsname\else
  \expandafter\newif\csname iflyraInglese\endcsname
\fi

% --- nomi del prodotto ------------------------------------------------------
\newcommand{\Lyra}{\textsc{Lyra}}
\newcommand{\LyraCapital}{\textsc{Lyra Capital}}
\newcommand{\app}{\iflyraInglese the application\else l'applicazione\fi}
\newcommand{\App}{\iflyraInglese The application\else L'applicazione\fi}
\newcommand{\sito}{\iflyraInglese the site\else il sito\fi}
\newcommand{\Sito}{\iflyraInglese The site\else Il sito\fi}

% --- i pacchetti del motore -------------------------------------------------
%  Si scrivono \pkgCore e non \file{core/}: così il giorno in cui una cartella
%  cambia nome, cambia qui.
\DeclareRobustCommand{\pkg}[1]{{\ttfamily\small\color{lyraInchiostro!78}#1}}
\newcommand{\pkgCore}{\pkg{core/}}
\newcommand{\pkgStrategie}{\pkg{strategies/}}
\newcommand{\pkgApprocci}{\pkg{approaches/}}
\newcommand{\pkgPesi}{\pkg{weights/}}
\newcommand{\pkgOverlay}{\pkg{overlays/}}
\newcommand{\pkgRegime}{\pkg{regime/}}
\newcommand{\pkgDati}{\pkg{data/}}
\newcommand{\pkgBacktest}{\pkg{backtest/}}
\newcommand{\pkgLive}{\pkg{live/}}
\newcommand{\pkgEsecuzione}{\pkg{execution/}}
\newcommand{\pkgReport}{\pkg{reporting/}}
\newcommand{\pkgDashboard}{\pkg{dashboard/}}
\newcommand{\pkgDominio}{\pkg{domain/}}
\newcommand{\pkgPorte}{\pkg{ports/}}
\newcommand{\pkgApplicazione}{\pkg{application/}}
\newcommand{\pkgAdattatori}{\pkg{adapters/}}
\newcommand{\pkgComposizione}{\pkg{composition/}}
\newcommand{\pkgApi}{\pkg{api/}}

% --- termini del dominio, con la prima occorrenza indicizzata ---------------
\newcommand{\term}[1]{\emph{#1}\index{#1}}
\newcommand{\termdef}[1]{\textbf{#1}\index{#1}}

% --- riferimenti a classi, funzioni, file -----------------------------------
%  \simbolo{Nome}            una classe o una funzione
%  \sorgente{file.py}{120}   un punto preciso del sorgente
\DeclareRobustCommand{\simbolo}[1]{\cod{#1}}
\newcommand{\sorgente}[2]{{\ttfamily\footnotesize\color{lyraInchiostro!60}#1:#2}}

% --- tasti e voci dell'interfaccia, per il manuale utente -------------------
%  \voce{Overview}           una voce di menu o una scheda
%  \pulsante{Run backtest}   un pulsante da premere
%  \campo{Target vol}        un campo da compilare
%  \percorso{Settings > Data} un percorso di navigazione
\newcommand{\voce}[1]{\textsf{\textbf{#1}}}
\newcommand{\pulsante}[1]{%
  \tcbox[on line, boxsep = 1pt, left = 3pt, right = 3pt, top = 0.5pt,
         bottom = 0.5pt, arc = 3pt, boxrule = 0.5pt,
         colback = lyraInchiostro, colframe = lyraInchiostro,
         coltext = white]{\sffamily\scriptsize\bfseries #1}}
\newcommand{\pulsantesec}[1]{%
  \tcbox[on line, boxsep = 1pt, left = 3pt, right = 3pt, top = 0.5pt,
         bottom = 0.5pt, arc = 3pt, boxrule = 0.5pt,
         colback = white, colframe = lyraInchiostro!50,
         coltext = lyraInchiostro]{\sffamily\scriptsize\bfseries #1}}
\newcommand{\campo}[1]{\textsf{\itshape #1}}
\newcommand{\percorso}[1]{\textsf{\small #1}}
\newcommand{\freccia}{\;\textcolor{lyraAccento}{$\rightarrow$}\;}

% --- una rotta HTTP ---------------------------------------------------------
%  \rotta{GET}{/api/marketplace/<id>/perche-non-mio}
%
%  Esiste perche' \cod{} e' una scatola indivisibile: in una colonna di
%  tabella larga cinque centimetri un percorso lungo sborda, e non c'e'
%  niente che TeX possa fare. Qui il percorso e' testo normale e puo' andare
%  a capo dopo ogni barra, che e' il punto naturale.
%
%  Il ciclo prende un carattere alla volta: funziona perche' un percorso non
%  contiene spazi (un argomento non delimitato li mangerebbe) ed e' in ASCII
%  (una lettera accentata e' un carattere attivo e manda in errore il ciclo).
%  Per qualunque altra cosa si usa \cod{}.
\makeatletter
\def\lyra@finePercorso{\lyra@finePercorsoMARK}
\def\lyra@percorso#1{%
  \ifx#1\lyra@finePercorso
    \let\lyra@prossimo\relax
  \else
    #1%
    % Dopo questi quattro caratteri la riga puo' andare a capo. Sono i punti
    % in cui un percorso o un nome di file si leggono comunque bene spezzati.
    \ifx#1/\allowbreak\else
      \ifx#1\_\allowbreak\else
        \ifx#1.\allowbreak\else
          \ifx#1-\allowbreak\fi\fi\fi\fi
    \let\lyra@prossimo\lyra@percorso
  \fi
  \lyra@prossimo}
\DeclareRobustCommand{\percorsoHTTP}[1]{%
  {\ttfamily\small\color{lyraInchiostro!85}%
   \lyra@percorso#1\lyra@finePercorso}}
\makeatother

\DeclareRobustCommand{\rotta}[2]{%
  {\sffamily\scriptsize\bfseries\color{lyraAccentoScuro}#1}\,\percorsoHTTP{#2}}

% --- la figura standard -----------------------------------------------------
%  Centra, mette la didascalia sotto e l'etichetta. Esiste perché le tre righe
%  scritte a mano ogni volta divergono: una figura senza \label non si può
%  citare, e una citata per numero si rompe appena se ne aggiunge un'altra.
%  #1 file da includere (senza estensione)  #2 didascalia  #3 etichetta
%  `max width` e non `width`: una figura piccola resta della sua misura, una
%  piu' larga della giustezza viene rimpicciolita invece di sbordare nel
%  margine. Senza, ogni figura andrebbe dimensionata a mano e ridimensionata
%  a ogni modifica del testo.
\newcommand{\figuraTikz}[3]{%
  \begin{figure}[htbp]
    \centering
    \begin{adjustbox}{max width = \linewidth}
      \input{figure/tikz/#1}
    \end{adjustbox}
    \caption{#2}\label{fig:#3}
  \end{figure}}

%  Come sopra ma per un'immagine, con larghezza.
%  #1 larghezza  #2 file  #3 didascalia  #4 etichetta
\newcommand{\figuraImg}[4]{%
  \begin{figure}[htbp]
    \centering
    \includegraphics[width = #1\linewidth]{#2}
    \caption{#3}\label{fig:#4}
  \end{figure}}

%  Schermata con cornice: una schermata senza bordo si confonde con il fondo
%  della pagina, e non si capisce dove finisce l'interfaccia. La cornice sta
%  DENTRO la larghezza chiesta: altrimenti una schermata larga quanto la
%  pagina sborda di un punto per lato.
\newcommand{\figuraSchermata}[4]{%
  \begin{figure}[htbp]
    \centering
    \setlength{\fboxsep}{0pt}\setlength{\fboxrule}{0.5pt}%
    {\color{lyraGrigioChiaro}\fbox{\includegraphics[width = \dimexpr #1\linewidth - 2\fboxrule\relax]{#2}}}
    \caption{#3}\label{fig:#4}
  \end{figure}}

% --- riferimenti incrociati con la parola davanti ---------------------------
%  In inglese la parola vuole la maiuscola («see Chapter 4»), in italiano no.
\newcommand{\vediFig}[1]{\iflyraInglese Figure\else figura\fi~\ref{fig:#1}}
\newcommand{\vediTab}[1]{\iflyraInglese Table\else tabella\fi~\ref{tab:#1}}
\newcommand{\vediCap}[1]{\iflyraInglese Chapter\else capitolo\fi~\ref{cap:#1}}
\newcommand{\vediSez}[1]{\iflyraInglese Section~\else \S\,\fi\ref{sez:#1}}

% --- numeri e unità ---------------------------------------------------------
%  Ogni lingua con la sua convenzione, la stessa dell'applicazione e del sito:
%  in italiano «10\,%» e «49\,€», in inglese «10%» e «€49» (le migliaia, in
%  inglese, si scrivono con la virgola nell'argomento: \eur{1,490}).
\newcommand{\perc}[1]{\iflyraInglese #1\%\else #1\,\%\fi}
\newcommand{\eur}[1]{\iflyraInglese \texteuro#1\else #1\,\texteuro\fi}


% ===========================================================================
%  Le macro dentro i segnalibri del PDF
% ===========================================================================
%  hyperref deve trasformare il titolo di ogni capitolo e sezione in una
%  stringa semplice per il pannello dei segnalibri. Le macro che compongono
%  qualcosa — una scatola colorata, un ciclo sui caratteri — lo mandano in
%  tilt: l'errore e' «TeX capacity exceeded [input stack size]» su una riga
%  che non ha niente di strano, tipo
%      \subsection{Che cosa c'e' sotto \file{data/}}
%
%  Qui si dice a hyperref di considerarle, nei soli segnalibri, come
%  «stampa il contenuto e basta».
\pdfstringdefDisableCommands{%
  \makeatletter
  \let\file\@firstofone
  \let\cod\@firstofone
  \let\pkg\@firstofone
  \let\simbolo\@firstofone
  \let\percorsoHTTP\@firstofone
  \let\voce\@firstofone
  \let\campo\@firstofone
  \let\pulsante\@firstofone
  \let\pulsantesec\@firstofone
  \let\percorso\@firstofone
  \let\term\@firstofone
  \let\termdef\@firstofone
  \def\rotta#1#2{#1 #2}%
  \def\freccia{ > }%
  \def\Lyra{Lyra}%
  \def\LyraCapital{Lyra Capital}%
  \makeatother
}

% ===========================================================================
%  Controllo: la sillabazione italiana c'è davvero?
% ===========================================================================
%  I pattern di sillabazione si caricano quando si GENERA il formato, non
%  quando si compila: se mancano, babel non fallisce — ripiega sull'inglese, e
%  il documento esce con le parole spezzate nei punti sbagliati. È un difetto
%  che si vede solo guardando la pagina, quindi vale la pena dirlo nel log.
%
%  Per rimediare, su una TeX Live completa: `tlmgr install hyphen-italian`;
%  su Debian/Ubuntu: `apt install texlive-lang-italian`; poi `fmtutil-sys
%  --byfmt pdflatex`.
%
%  Il controllo guarda \language, che è un parametro intero del motore e si
%  può sempre confrontare: 0 è la lingua di partenza (inglese). Se dopo aver
%  selezionato l'italiano siamo ancora lì, i pattern non ci sono.
\AtBeginDocument{%
  \ifnum\language = 0
    \PackageWarningNoLine{lyra}{%
      Pattern di sillabazione ITALIANI assenti: il testo sara' spezzato\MessageBreak
      con le regole inglesi. Installare hyphen-italian e rigenerare il\MessageBreak
      formato (fmtutil-sys --byfmt pdflatex)}%
  \fi}
